
\section{Objetivo}
\label{inicio_cap5}

El control documental en Skyfilms asegura que todos los documentos, registros, instructivos, certificados y evidencias relacionadas con la operación UAS estén actualizados, disponibles y bajo custodia segura. Este control se implementa conforme a los lineamientos de la Aerocivil y hace parte integral del Sistema de Gestión de la Seguridad Operacional (SMS).\\

Las Especificaciones Relativas a las Operaciones (OpSpecs) forman parte integral del MO y se controlan como Anexo con versión y tabla de revisiones independientes.


\section{Control Documental}

\subsection{Documentos Controlados}
El sistema de control documental de Skyfilms cubre, como mínimo:
\begin{itemize}
    \item Manual de Operaciones (MO) y sus revisiones.
    \item Manual de SMS y documentos asociados.
    \item Manual Control de Mantenimiento (MCM).
    \item Instructivos, procedimientos y formatos operacionales.
    \item Hojas de vida y certificados de idoneidad del personal UAS (pilotos, observadores, Jefe de Pilotos, gerente SMS).
    \item Autorizaciones de vuelo emitidas por la Aerocivil.
    \item Registros de vuelo (libro de vuelo, bitácoras de pilotos).
    \item Grabaciones de video, audio y registros electrónicos (logs de vuelo).
    \item Certificados de registro de UAS y equipos tecnológicos asociados.
\end{itemize}

\subsection{Responsables del Control Documental}
\begin{itemize}
    \item El \textbf{Gerente SMS} será responsable del control y actualización de los manuales, instructivos y registros de seguridad operacional.
    \item El \textbf{Jefe de Pilotos} será responsable de la custodia y actualización de bitácoras de vuelo, hojas de vida de pilotos y observadores, así como de la documentación técnica asociada a cada UAS.
\end{itemize}


\subsection{Conservación}
Los documentos y registros serán almacenados en:
\begin{enumerate}
    \item \textbf{Repositorio digital seguro}: Carpeta en la nube corporativa de Skyfilms con acceso restringido al Gerente General, Gerente SMS y Jefe de Pilotos. Incluye copias de seguridad automáticas semanales.
    \item \textbf{Repositorio físico}: Carpeta archivadora bajo custodia del Gerente SMS, destinada a los documentos que requieran soporte físico o firma original.
\end{enumerate}



\subsection{Control de Versiones}
\begin{enumerate}
    \item Toda modificación de manuales, instructivos o formatos será propuesta por el área correspondiente y aprobada por el Gerente General.
    \item El Gerente de Seguridad Operacional registrará la enmienda en el historial de cambios, informando a todo el personal involucrado mediante circular interna.
    \item Se conservarán versiones anteriores durante un periodo mínimo de \textbf{cinco (5) años} para efectos de trazabilidad.
\end{enumerate}





\section{Control de Registros}

Skyfilms S.A.S. mantiene un sistema de registros operacionales que permite demostrar la planificación, ejecución, supervisión y cierre de todas las operaciones desarrolladas por la organización.

Los registros operacionales constituyen evidencia objetiva del cumplimiento de los procedimientos establecidos en el presente Manual de Operaciones y forman parte del Sistema de Gestión Documental de Skyfilms.


\subsection{Objetivo}

Los registros operacionales tienen como finalidad:

\begin{itemize}

\item Garantizar la trazabilidad de las operaciones.

\item Demostrar el cumplimiento de los procedimientos establecidos.

\item Facilitar auditorías internas y externas.

\item Apoyar la investigación de incidentes.

\item Generar información para la mejora continua.

\end{itemize}


\renewcommand{\arraystretch}{1.4}

\begin{longtable}{|p{4cm}|p{5cm}|p{3cm}|p{2cm}|}
\hline

\textbf{Registro} &
\textbf{Responsable} &
\textbf{Conservación} &
\textbf{Código}

\\ \hline

Libro de Vuelo &
Piloto Remoto &
Según MDM &
F-001

\\ \hline

Checklist Preoperacional &
Piloto Remoto &
Según MDM &
F-002

\\ \hline

Checklist Postoperacional &
Piloto Remoto &
Según MDM &
F-003

\\ \hline

Registro de Baterías &
Piloto Remoto / Mantenimiento &
Según MDM &
F-004

\\ \hline

Reporte de Mantenimiento y Orden de Trabajo &
Piloto Remoto / Técnico &
Según MDM &
F-005

\\ \hline

Hoja de Vida de la Aeronave &
Jefe de Pilotos &
Según MDM &
F-007

\\ \hline

Solicitud de Servicio &
Cliente / Jefe de Pilotos &
Según MDM &
F-008

\\ \hline

Planificación Operacional &
Jefe de Pilotos &
Según MDM &
F-009

\\ \hline

Briefing &
Piloto Remoto &
Según MDM &
F-010

\\ \hline

Evaluación del Riesgo &
Jefe de Pilotos &
Según MDM &
F-011

\\ \hline

Reporte SMS &
Todo el personal &
Según SMS &
F-006

\\ \hline

Reporte de Peligros &
Gerente SMS &
Según SMS &
F-012

\\ \hline

Reporte de Incidente o Accidente &
Piloto Remoto / Gerente SMS &
Según SMS &
F-013

\\ \hline

\end{longtable}


\subsection{Conservación de registros}

Todos los registros operacionales deberán mantenerse bajo control documental conforme a lo establecido en el Capítulo 5 del presente Manual y en la Matriz Documental Maestra (MDM).

Los formatos de estos registros se encuentran en el Capítulo 6 de este Manual.

Los registros podrán conservarse en formato físico, digital o electrónico, siempre que se garantice su integridad, disponibilidad, trazabilidad y protección contra modificaciones no autorizadas.


\subsection{Sistema de gestión documental}
Skyfilms administrará sus registros operacionales mediante un Sistema de Gestión Documental centralizado en la nube, el cual permitirá controlar la creación, revisión, almacenamiento, consulta y disposición final de la información. 

Los formatos operacionales, listas de verificación (checklists) y reportes de riesgo se gestionarán a través de formularios digitales estandarizados o bitácoras físicas controladas, garantizando en todo momento la integridad, trazabilidad y protección de los datos contra modificaciones no autorizadas, manteniéndolos siempre a disposición de la Autoridad Aeronáutica.

\subsection{Trazabilidad operacional}

Toda operación desarrollada por Skyfilms deberá permitir la trazabilidad completa de la misión desde la recepción del servicio hasta el cierre documental.

Como mínimo deberá ser posible identificar:

\begin{itemize}

\item Cliente.

\item Fecha.

\item Aeronave utilizada.

\item Piloto Remoto.

\item Observador Visual.

\item Lugar de operación.

\item Evaluación del riesgo.

\item Registros de vuelo.

\item Estado técnico de la aeronave.

\item Reportes SMS asociados.

\end{itemize}

\begin{tcolorbox}[title=Documentos asociados,colback=gray!5,colframe=black]

\textbf{Procedimiento relacionado}

PO-007 Gestión de Registros Operacionales.

\vspace{0.3cm}

\textbf{Documentos relacionados}

\begin{itemize}

\item Matriz Documental Maestra (MDM)

\item Todos los formatos operacionales

\item Manual SMS

\item Manual MCM

\end{itemize}

\end{tcolorbox}


